\documentclass[11pt,letterpaper]{article}
\usepackage[T1]{fontenc}
\usepackage{lmodern,amsmath,amssymb,amsthm,microtype}
\usepackage[margin=1in]{geometry}
\usepackage[hidelinks]{hyperref}
\providecommand{\tightlist}{\setlength{\itemsep}{0pt}\setlength{\parskip}{0pt}}
\DeclareUnicodeCharacter{2264}{\ensuremath{\le}}
\DeclareUnicodeCharacter{2265}{\ensuremath{\ge}}
\DeclareUnicodeCharacter{2260}{\ensuremath{\ne}}
\DeclareUnicodeCharacter{2212}{\ensuremath{-}}

\DeclareUnicodeCharacter{2211}{\ensuremath{\sum}}

\usepackage{longtable,booktabs,array}
\title{The eighth weight is at least 121 when the largest is between 149 and 150}
\author{Aafi Khan Malek\\ \small Independent researcher, Ahmedabad, India\\
\small\href{mailto:aafimalek2032@gmail.com}{aafimalek2032@gmail.com}}
\date{28 November 2026}
\begin{document}
\maketitle
\textbf{Theorem.} Let nine sorted positive real weights have pairwise
unit-separated subset sums. If \(149\le x_9<150\), then \(x_8\ge121\).
This excludes a slice of the real optimization problem, not the whole
range below 150. We use the
\href{https://github.com/Aafimalek/novel-discoveries/tree/4c5e61d8f6b55fa4fc9cf31aff871a1c0a6c4cc2/nine-separated-subset-sums/verification/rank-cuts-below-150/hall-rank-runs}{rank-run argument}
from the previously certified eight Harper cuts and forbidden rank-149
edges. Exact integer and rational scripts replay the finite box
exclusions; the new theorem has not yet been formalized in Lean.

\section{A cap forces an entire rank-edge
run}\label{a-cap-forces-an-entire-rank-edge-run}

Let \(t\) be any integer bound for \(\lfloor x_8\rfloor\). At each
Harper cut \(k\) in 158,159,161,162,168,169,171,172, every rank
\(k+1,\dots,k+148\) is an external neighbor, by the established
forbidden-band saturation. For \(j\) with
\(\max(k+t+1,257)\le j\le k+148\), the edge from before the cut to \(j\)
must change coordinate 9: its rank length exceeds \(t\). Its length is
at most 148, because a length-149 edge would begin at \(j-149\) in the
forbidden rank band. The lower ranks obtained by subtracting \(M\)
strictly increase with \(j\), are at least \(j-148\), and end at \(k\).
Consequently each lower rank is exactly \(j-148\). This forces every
start from \(\max(109,k+t-147)\) through \(k\). The unions for
\(t\le119\) considered here are single runs, as the replay checks. Their
associated full-sign slots \(q=k-108\) run from
\(\max(109,158+t-147)-108\) to 64.

\section{Strengthened rank-slot
moment}\label{strengthened-rank-slot-moment}

Write \(u_1<\cdots<u_{128}\) for the absolute eight-weight full-sign
values (complementary signs counted once), and \(N(v)=\#\{i:u_i<v\}\).
The \(u_i\) and the corresponding nine-weight magnitudes \(d_i\) each
have consecutive spacing at least 2. A forced start with slot \(q\)
gives \(\alpha_q=u_{j_q}\) with \(149-2q\le\alpha_q<151-2q\) and
full-sign positions \(q\) and \(149-q\) at \(M-\alpha_q\) and
\(M+\alpha_q\). Set \(a_q=N(2M-\alpha_q)\) and \(b_q=N(2M+\alpha_q)\).
Ranking the two full forms gives \(a_q+b_q=148\) and \(j_q=a_q-q+1\).
Since the \(\alpha_q\) strictly decrease, the \(a_q\) are nondecreasing
but \(a_{q+1}-a_q\le0\) for consecutive slots. Thus all \(a_q=a\) along
the run, with \(64\le a\le74\).

The next \(u_{a+1}\) is strictly above \(2M-\alpha_q\), and the full
form \(u_{a+1}-M\) must lie at least 2 above \(M-\alpha_q\). The same
reasoning at \(2M+\alpha_q\) compares to \(M+\alpha_q\). Therefore the
following are valid point bounds:

\begin{verbatim}
u[a-q+1] ≥ 149−2q
u[a+1]   ≥ 150+2q
u[149-a] ≥ 449−2q.
\end{verbatim}

Each bound advances by at least 2 per rank; also \(u_j\ge2j-1\). The
exact profiles for all \(64\le a\le74\) give the moment
\(\sum_{i=1}^8 x_i^2=\sum_j u_j^2/128\) by sign averaging. Only lower
envelopes are used; no assumption that such a profile is attainable.

\section{Excluding caps through 118}\label{excluding-caps-through-118}

Write \(y=(t-6,\ldots,t+1)\), so \(x_i\le y_i\) for \(i\le8\). Put
\(r_i=y_i-x_i\) and \(R=\sum r_i\). The deficits are nonincreasing and
nonnegative by singleton spacing. If \(B\) is the moment lower bound,
the squared loss \(D=\sum y_i^2-B\) is an upper bound for
\(\sum_i (2y_i r_i-r_i^2)\). Negative \(D\) immediately excludes the
cap; this works for the full \(x_8<110\) shard at \(t=109\), and for
\(t=110\). For the later caps \(R\le D/y_1\) because \(r_i\le y_i\).

At \(y\) the balanced signs (four plus, four minus) have magnitudes at
most 16. All other signs have magnitude at least \(2t-20\), achieved by
taking the smallest five positive. Replacing \(y\) by \(x\) moves every
magnitude by at most \(R\). For caps 111--119 the forced
\(\alpha_{64}<23\) cannot come from an unbalanced sign, by the coarse
\(R\) bound. Hence for some balanced sign
\(|\sum_i \sigma_i r_i|\le39\). As each \(r_i\) lies on one of the two
sides, \(r_i\le (R+39)/2\). For \(R\ge39\) this gives
\(\sum_i r_i^2\le (R^2+39^2)/2\): if the largest entry is at most
\(R/2\), use \(\sum r_i^2\le R^2/2\); otherwise concentrate the
remaining mass into one entry and increase the largest to its upper
bound. For \(R\le39\) simply use \(\sum r_i^2\le R^2\). Consequently
\(D\ge 2y_1R-\sum r_i^2\) gives the exact integer cutoffs used by the
replay. Both pieces are increasing on the relevant intervals (the low
piece starts with \(y_1>39\) and the high piece ends below \(2y_1\)).

There are 35 balanced complementary sign pairs. At slot \(q\), at least
\(a-q+1\ge65-q\) base magnitudes are below \(151-2q\). A cap admits at
most 35 plus the number of unbalanced \(y\) magnitudes below
\(151-2q+R\) in that region. This exact count excludes every cap from
111 through 118. For cap 119 it gives \(R<144\) but not a contradiction.

\section{Exact two-deficit boxes for cap
119}\label{exact-two-deficit-boxes-for-cap-119}

In this cap there are 43 distinct required \(\alpha\) slots. At least 8
have unbalanced sign. If \(R\le125\), an unbalanced form has magnitude
at least \(218-125=93\). Only the seven slots \(q=22,\ldots,28\) can
reach 93 (slot 29 has strict upper 93). Thus \(R>125\); the box replay
relaxes to \(R\ge125\).

The deficits \(r_1,r_2\) lie in \([0,112]\). Cover them by unit cells
\([A,A+1]\times[B,B+1]\) with integer \(B\le A\). This covers all
ordered pairs, including integral boundaries. Let \(T=r_3+\cdots+r_8\).
Taking the lower endpoints \(r_1=A\) and \(r_2=B\), the first two
minimum squared losses are \(2y_1A-A^2\) and \(2y_2B-B^2\), since
squared loss is increasing for \(r_i\le y_i\). Each tail entry is at
most \(B+1\), so its loss is at least \((2y_3-B-1)r_i\). This gives the
rational upper bound for \(T\). The constraints \(125\le R\le144\) and
\(0\le T\le6(B+1)\) give the other simple bounds. All upper interval
rounding is outward.

For a fixed sign pattern, inverse-prefix differences express the sorted
tail deficits as a nonnegative linear combination of initial-segment
indicators. For each indicator take the average of its signs.
Consequently the signed contribution of the tail lies between \(T\)
times the least and greatest of its six prefix-sign averages. All these
averages have denominator dividing 60. Combining with the first two
deficit intervals and \(y\) bounds each full-sign magnitude by an
interval. Its eligible \(q\) slots form an integer interval, since
\(\alpha_q\in[149-2q,151-2q)\). Every required slot must be matched to a
different one of the 128 complementary full sign pairs; there is no
reliance on simultaneous realizability of this outer approximation.

The replay attempts exact interval matching, in ascending \(q\),
choosing an eligible interval with smallest right endpoint. The greedy
rule is complete by exchange: if an assignment uses a different interval
at the current \(q\), move the earliest-deadline interval to \(q\) and
put the displaced one at its previous later slot, if any. This is valid
because the displaced interval contains the current \(q\) and ends no
earlier than the chosen one. All 6,328 cells are excluded: 1,597 by
moment loss, 2,631 by tail bounds, and 2,100 by interval matching. No
heavy solver, floating comparison, or proof kernel is used for this
replay.

\section{The final cap of 120}\label{the-final-cap-of-120}

Here \(M=x_9\), the cap is \(t=120\), and the two-deficit moment uses
the same notation as the cap-119 proof. The consecutive forced run is
\(q=23,\ldots,64\). In particular each \(\alpha_q\) is one unsigned base
eight-weight full sign value, at exact index \(j_q=a-q+1\). The common
counts \(a,b\) are as before, with \(64\le a\le74\) and \(b=148-a\).

The \(u\)-values before index \(j_{64}\) are all strictly below 21,
because \(\alpha_{64}<23\) and consecutive \(u\) spacing is at least 2.
Values after \(j_{23}\) are at least 105 because \(\alpha_{23}\ge103\).
Thus exactly \(a-64\) values are below 21, and one is in each half-open
interval \([149-2q,151-2q)\) for \(q=64,63,\ldots,23\).

The previous packet uses a +2 lower sharpening of threshold ranks. There
is an upper sharpening too. The full nine magnitude \(M-\alpha_q\)
cannot coincide with \(|M-u_i|\) for \(u_i\) near \(2M-\alpha_q\) (nor
\(M+\alpha_q\) near \(2M+\alpha_q\)); every different full magnitude is
at least 2 away from it. Therefore

\begin{verbatim}
u[a]    ≤ 2M−alpha[q]−2
u[a+1]  ≥ 2M−alpha[q]+2
u[b]    ≤ 2M+alpha[q]−2
u[b+1]  ≥ 2M+alpha[q]+2.
\end{verbatim}

For the first upper bound, if \(u_a\le M\), then it already satisfies
the bound because \(M-\alpha_q=d_q\ge45\). Otherwise compare the
magnitudes \(u_a-M\) and \(M-\alpha_q\). The second upper bound is
analogous. No ties at the threshold are possible. With the extreme
\(q\)-values and \(149\le M<150\) this gives \(u_a<195\),
\(u_{a+1}\ge278\), \(u_b<321\), \(u_{b+1}\ge403\).

Hence the allowed bins for all 128 complementary base sign pairs, with
exact capacities, are:

\begin{longtable}[]{@{}lr@{}}
\toprule\noalign{}
Half-open range & Number of values \\
\midrule\noalign{}
\endhead
\bottomrule\noalign{}
\endlastfoot
{[}1,21) & \(a-64\) \\
each {[}149−2q,151−2q), q=64 down to 23 & 1 \\
{[}105,195) & 22 \\
{[}278,321) & \(148-2a\) \\
{[}403,1000) & \(a-20\) \\
\end{longtable}

The last upper endpoint 1000 is harmless, since all eight weights are at
most 121. Intervals omitted from this table are forbidden. The
zero-capacity possibility \(a=b=74\) is supported.

\section{Exact two-stage proof}\label{exact-two-stage-proof}

The 42 forced \(\alpha\)'s require at least seven unbalanced signs, but
slots 23 through 28 are the only ones whose strict upper bounds exceed
93. The minimum unbalanced magnitude at the cap vector
\(y=(114,\ldots,121)\) is 220. As each signed magnitude can move by at
most the total deficit \(R\), this gives \(R>127\), safely relaxed to
\(R\ge127\) in replay.

\begin{enumerate}
\def\labelenumi{\arabic{enumi}.}
\tightlist
\item
  Cover \(r_1,r_2\) by integer unit boxes and bound remaining deficit by
  the moment loss, as in the cap-119 packet. Using the minimum moment
  across all \(a\) gives 291 cells that have not yet been excluded by
  the 42-slot-only matching.
\item
  Split each of these cells into the 11 possible counts
  \(a=64,\ldots,74\). Use its own moment bound, stronger deficit bound,
  and the finite bins above. This is 3,201 cell/count profiles.
  Loss/tail bounds reject 2,498, and exact interval matching rejects
  every one of the other 703.
\end{enumerate}

For the strengthened total-deficit bound in step 2, all \(r_i\le113\).
The coarse bound is at most \(D/114<226\) (where \(D\) is the profile's
moment loss budget). If \(R>113\), the inequality
\(\sum r_i^2\le113^2+(R-113)^2\) follows by concentrating the remaining
mass, or by the same largest-entry argument as in the cap-119 proof.
Thus \(2\cdot114 R-\sum r_i^2\) gives a monotone lower bound on loss and
yields the integer cutoffs in replay. If the resulting \(R\) is at most
\(220-23\), then \(\alpha_{64}<23\) requires a balanced sign and the
separate 39-bound R refinement from the cap-119 proof also applies. None
of these steps presumes its relaxation is attainable.

For each profile use the outward-rounded rational interval of possible
magnitudes for each sign pair. In the ordered list of 128 unit-capacity
positions obtained by expanding the bins, the positions intersecting
that interval are consecutive. Run earliest-deadline matching against
all 128 positions; the same exchange argument as in the cap-119 proof
proves greedy completeness even for repeated bin boundaries. If there is
a true weight vector, each sign pair would map to a distinct position,
so failure proves exclusion.

\end{document}
